\documentclass[11pt]{article}
\usepackage[UTF8]{ctex}
\usepackage{geometry}
\geometry{a4paper, top=2.5cm, bottom=2.5cm, left=2.5cm, right=2.5cm}
\usepackage{amsmath}
\usepackage{amssymb}
\usepackage{amsthm}
\usepackage{booktabs}
\usepackage{tabularx}
\usepackage{longtable}
\usepackage{enumitem}
\usepackage{xcolor}
\usepackage{tcolorbox}
\usepackage{algorithm}
\usepackage{algpseudocode}
\usepackage[colorlinks=true, linkcolor=blue!60!black, citecolor=darkgray, urlcolor=blue!60!black, bookmarks=true, bookmarksnumbered=true, unicode=true]{hyperref}
\widowpenalty=10000
\clubpenalty=10000

\newtheorem{theorem}{定理}
\newtheorem{lemma}[theorem]{引理}
\renewcommand{\proofname}{证明}

\newtcolorbox{keybox}[1][]{colback=blue!4, colframe=blue!40!black, boxrule=0.6pt, arc=2pt, left=6pt, right=6pt, top=4pt, bottom=4pt, #1}

\title{空}
\date{}

\begin{document}
\tableofcontents
\newpage

\section{目的与完成度声明}

本资料整合本工作的\textbf{算法理论}与\textbf{实验佐证}，供同态密码专家进行双重审核。全部数字可溯源至代码仓库 \texttt{repro/} 工件与 git 提交记录。

\begin{keybox}
\textbf{完成度（截至本资料生成）}：安全组合定案、终版性能主表、层安全审计、容量曲线（h=41/$\sigma$+11 口径）、TFHE-rs 同机锚点、occupancy/break-even、DFR 三层证据、多参数集修正推导（4096/5-bit 已\textbf{实测}）均已\textbf{完成}。本轮 stage368 收官新增：r=8 静机复测、E-PF 融合 A/B、postops 交替翻转分诊（5 变体，判为测量伪象）、S5 多参数集实测（4096 门测试 0/4096 失配）。在途（不改变主表方向）：10-trial 噪声@最终参数、E4 新估计器独立复算臂（V5A2）。\textbf{勘误}：早前 r-lane 表中路径A $4.32$ ms/msg 系误取 pathA 构建二进制的 stock 对照列（跨构建代码布局漂移实测达 $23\%$）；本版全部 r-lane 结论已改为\textbf{同构建内部对照}口径（表~\ref{tab:rlane}）。
\end{keybox}

\section{参数体系（$\theta$ 向量，实码值）}

表~\ref{tab:theta} 给出最终参数（Final，SET\_2\_3\_2048 族）的完整参数向量。所有密钥层的安全含义在 \S5 逐一审计。

\begin{table}[htbp]
\centering
\caption{最终参数向量（代码可验证值）}
\label{tab:theta}
\small
\begin{tabular*}{\columnwidth}{@{\extracolsep{\fill}} l l l}
\toprule
符号 & 值 & 说明 \\
\midrule
$n=N_{\mathrm{acc}}=N_{\mathrm{out}}$ & 2048 & 输入/累加器/输出环维数 \\
$q$（torus） & $2^{64}$ & \\
输入密钥 $s_{\mathrm{in}}$ & 二值稀疏 $h{=}42$，条件于最大间隙 $\le 2^7{-}1$ & 拒绝采样 \\
$\sigma_{\mathrm{in}}$ & $2^{-15}$（绝对 $2^{49}$） & \\
输出密钥 $s_{\mathrm{out}}$ & 三元 $h_{\mathrm{out}}{=}512$，\textbf{交替符号} & 如实建模 \\
$\sigma_G$（BSK） & $2^{-50}$（原值，\textbf{不修正}） & 见 \S5.1 \\
packing 密钥 & 二值 $h{=}256$，$\sigma_{\mathrm{PKS}}{=}2^{-44}$ & \\
$Q$（SQ 尺度） & $2^{16}$ & \S4 \\
gadget（stock/路径A） & $B_g{=}2^{23}$，$\ell{=}1$ & \\
$\rho$（$r\_prec$） & 7 & \\
消息精度 $p$ & 3（族内 5/7/9 同 $h{=}42$） & \\
\bottomrule
\end{tabular*}
\end{table}

\section{三个自举算法}

\subsection{686 基线（标量，已发表）}

输入 TRLWE $\mathrm{in}=(a,b)$（$s_{\mathrm{in}}$ 下）与测试向量 TV。初始化 $acc_i = X^{b_i}\!\cdot\!TV$（$i\in[n)$）。对 $s_{\mathrm{in}}$ 的 $h{+}1$ 个稀疏步，每步对\textbf{整个 $n$ 数组}做单项乘（每精度位：换位槽 NCMUX = $X^{-1}$ 自同构 + KS 后 CMUX；直达槽 CMUX）并做数位减 SubA。外积总数 $S=(h{+}1)\rho n$。尾声为逐槽抽取 + packing KS + hw-KS。

\subsection{矩阵外积拓展（本文贡献 1；路径A为其工程优化）}

累加器元素由 TRLWE 改为 \textbf{PVW-TMLWE}（1 个掩码 $a$ 与 $r$ 个体 $b_j$ 共享），选择子由 TRGSW 改为 \textbf{MAT-TRGSW}（$(1{+}r)\ell$ 行）。矩阵外积
\[
\mathrm{MV\text{-}EP}(c,C)\;=\;\sum_{i\le r,\; j<\ell}\ \mathrm{Decomp}_j(c_i)\otimes C[i\ell+j]
\]
一次盲旋转服务 $n\times r$ 个消息--LUT 查询；外积计数不变而每 lane 摊销 $\div r$；分解行数由 $2r$ 降至 $(1{+}r)$（共享掩码的结构性节省）。

\begin{algorithm}[htbp]
\caption{矩阵形式摊销稀疏自举（外层）}
\label{alg:mat}
\begin{algorithmic}[1]
\Require $\mathrm{in}=(a,b)$；PVW 测试向量 $\mathit{tv}$（$r$ 路 LUT）；密钥 $\{C[\text{step}][\text{bit}]\}$、aut$_{-1}$、packing KS $\times r$、hw-KS
\Ensure $out[0..r{-}1]$
\For{$i\in[0,n)$} $acc_i \leftarrow X^{\lceil 2^{64} b_i\rceil_{2N}}\cdot \mathit{tv}$ \EndFor
\For{$\text{step}\in[0,h]$}
  \State $acc \leftarrow \mathrm{MonomialMul}(acc,\,C[\text{step}])$ \Comment{算法~\ref{alg:bfly}}
  \State $\mathrm{SubA}(acc,\,\bar a)$ \Comment{零数位快路径跳过整次 MV-EP}
\EndFor
\For{$\text{lane}\in[0,r)$}
  \State $ext \leftarrow$ 逐槽常数抽取$(acc,\text{lane})$
  \State $out[\text{lane}] \leftarrow \mathrm{KS}_{\mathrm{hw}}(\mathrm{PackingKS}(ext))$
\EndFor
\end{algorithmic}
\end{algorithm}

\begin{algorithm}[htbp]
\caption{MonomialMul：累加器数组的蝶形旋转（每精度位）}
\label{alg:bfly}
\begin{algorithmic}[1]
\For{$\text{bit}\in[0,\rho)$}，$w\leftarrow 2^{\text{bit}}$
  \For{$j\in[0,w)$} \Comment{换位槽：负旋转}
     \State $rot \leftarrow X^{-1}\!\cdot acc^{\mathrm{in}}_{n-w+j}$（自同构 + aut$_{-1}$ KS）
     \State $(out_j,\ out_{j+w}) \leftarrow \mathrm{DualSubCMUX}(acc^{\mathrm{in}}_j,\ rot,\ acc^{\mathrm{in}}_{j+w},\ C[\text{bit}])$
  \EndFor
  \For{$j\in[w,\,n{-}w)$} \Comment{直达槽}
     \State $out_{j+w} \leftarrow acc^{\mathrm{in}}_{j+w} + \mathrm{MV\text{-}EP}(C,\ acc^{\mathrm{in}}_j - acc^{\mathrm{in}}_{j+w})$
  \EndFor
\EndFor
\State 交换缓冲
\end{algorithmic}
\end{algorithm}

DualSubCMUX 中一对输出共享一次 $X^{-1}$ 旋转与一次向量化双减：$(rot{-}p,\ rot{-}q)$ 一次算得，两输出各做一次 MV-EP。蝶形拓扑保证每个旋转尾恰参与 2 个输出的代数恒等式，故 $k{=}2$ 是结构上限（$k$-sub 推广被拓扑排除）。

\subsection{SQ 尺度量化外积（本文贡献 2）}

累加器全程维护在 $Q=2^{q}$（$q{=}16$）尺度：setup 量化 $\times 2^q$，每次外积后融合重缩放 $\gg (64{-}q)$ 保持尺度；选择子按 $B_g{=}2^{q}$ 加密，从而\textbf{gadget 分解被消除}（$Q$ 尺度累加器本身即单一数位）。NCMUX 为精确升移 $\ll(64{-}q)$、复用 stock 自同构 KS、舍入降移（每步一次 $q$-bit 舍入）。速度收益集中在标量（SQ/686 $=0.897$），随 $r$ 按 $\approx 1/r$ 衰减至平价；硬化收益双形态保持（\S4.3）。

\section{SQ 理论：正确性与噪声}

\begin{theorem}[消息路径精确性]
\label{thm:1}
在无回绕条件 $|c|_\infty \le 2^{q-1}$ 下，对盲旋转调度的每一步，
\[
\bigl(\mu\cdot 2^{64-q}\ \star\ c\bigr)\ \gg\ (64-q)\;=\;\mu\cdot c,
\]
且 SQ 盲旋转的输出消息与 stock 盲旋转（同 TV、同密钥）一致。
\end{theorem}

\begin{proof}
对调度步归纳。基础：setup 量化 $c_0=\mathrm{round}(\text{torus}\cdot 2^{q})$，消息占据高位，$|c_0|_\infty<2^q$，条件成立。归纳步：外积核为整数卷积（spqlios 低字对 $2^{64}$ 精确）加移位；整数卷积对消息线性，且无回绕条件保证中间乘积不溢出，故消息项经 $\times 2^{64-q}$ 再右移 $64{-}q$ 恰好还原为 $\mu\cdot c$；噪声项由引理~\ref{lem:1} 刻画。融合重缩放的舍入 $R$ 满足 $|R|_\infty\le 2^{q-1}$（单次 $q$-bit 舍入），维持尺度不变量与无回绕条件。归纳终止时累加器仍在 $Q$ 尺度，末尾 torus 提升为精确移位。故消息路径总舍入为零，仅有每次 NCMUX 一位的 $q$-bit 舍入噪声累积。
\end{proof}

\begin{lemma}[$\Delta$ 抑制]
\label{lem:1}
选择子噪声 $e$（标准差 $\sigma_G$）进入累加器的系数为：stock $\sigma\cdot 2^{27.1}$（KS 噪声 $\propto\sigma\cdot 2^{B_g}$，约 140 次/槽），SQ 为 $\sigma\cdot 2^{q+3.5}$；二者之比为
\[
\Delta \;=\; Q^2/T \;=\; 2^{\,2q-64}\;=\;2^{-32}\quad(q=16),
\]
即 $q$ 每降低 1 bit 换取 2 bit 的 $\sigma$ 余量。
\end{lemma}

推导要点：$Q$ 尺度下选择子噪声需经（$\times 2^{64-q}$ 谱乘）与（$\gg 64{-}q$ 重缩放）对消，净系数 $2^{-(64-q)}\cdot 2^{q}=2^{2q-64}$；而 stock 的分解数位放大 $B_g/2\approx 2^{22}$ 与约 140 次 KS 累积给出 $2^{27.1}$。

\textbf{理论优势三线（均有实测）}：(i) 速度——消分解使标量 SQ/686 $=0.897$（三配置 $0.894$--$0.898$ 稳定）；(ii) 余量——$\sigma{+}15$ 后 SQ 噪声 59.36，与引理预测 $59.1$ 相差 $<1$ bit，三档 $\sigma$ 均如此；(iii) 结构解耦——SQ 对 $\sigma_G$ 不敏感（平衡档实验：stock 付 $17\%$ 得 $+0.9$ bit，SQ 零收益零损失）。

\section{安全分析与参数定案}

\subsection{攻击组合与最终参数}

组合取各攻击之\textbf{最小值}（不叠加）。攻击面：uSVP、dual-hybrid、HD（合并论文公开代码）、lattice-MITM-hybrid（新版 lattice-estimator）、组合层闭式界 $T3'$（本工作）。

\begin{table}[htbp]
\centering
\caption{安全组合表（bit；$T3'$ 为本工作含条件熵修正的闭式组合界）}
\label{tab:sec}
\small
\begin{tabular*}{\columnwidth}{@{\extracolsep{\fill}} l ccccc c}
\toprule
层/参数 & uSVP & dual & HD & MITM-hyb & $T3'$ & min \\
\midrule
输入 $h{=}39$ & 898.6 & — & 不可行 & 126.2 & 124.5 & \textbf{124.5} \\
输入 $h{=}41$ & 902.2 & — & — & 130.0 & 130.5 & 130.0 \\
\textbf{输入 $h{=}42$} & 904.1 & — & — & \textbf{131.8} & 133.5 & \textbf{131.8} \\
输入 $h{=}52$ & 919.6 & — & — & 149.4 & 165.4 & 149.4 \\
\textbf{BSK $\sigma{=}2^{-50}$} & 211.6 & 318.3 & \textbf{136.2} & — & — & \textbf{136.2} \\
BSK $\sigma{=}2^{-39}$ & 310.5 & 229.6 & 175.1 & — & — & 175.1 \\
\bottomrule
\end{tabular*}
\end{table}

\begin{keybox}
\textbf{最终参数：$h{=}42$，$t{=}7$，$\sigma_G{=}2^{-50}$（原值）。}系统安全 $=\min(131.8,\,136.2)=131.8\ge 128$。三方独立验证：闭式 $T3'{=}133.5$、估计器 lattice-MITM $131.8$（两线相差 $<2$ bit 互证）、合并论文自身代码在 BSK 层 $136.2$。$h{=}39$（686 原参数）在正确条件熵建模下为 $124.5$，\textbf{不达标}。
\end{keybox}

\subsection{条件熵修正（本工作方法）}

实现中输入密钥经拒绝采样直至最大环形间隙 $\le 2^{t}-1$。均匀丢弃使条件分布均匀于接受族，故熵损失恰为 $-\log_2 p_{\mathrm{accept}}$；$p_{\mathrm{accept}}$ 由逐行复刻 \texttt{get\_min\_prec} 扫描语义的 Monte Carlo（$2\times10^4$ 样本）计算。修正后
\[
T3'(n,h,t)=\bigl(\log_2\tbinom{n}{h}\;-\;\mathrm{loss}(n,h,t)\;-\;\log_2 n\;-\;\delta\bigr)/2,\quad \delta=7.38 .
\]
实现的三元密钥采用\textbf{交替符号}（符号熵 1 bit 而非 $h$ bit），已如实计入建模；该层组合安全仍远高于格层，绑定层不变。

\subsection{层安全审计（最终参数下全层达标）}

\begin{table}[htbp]
\centering
\caption{逐层审计（min，bit）}
\label{tab:layers}
\small
\begin{tabular*}{\columnwidth}{@{\extracolsep{\fill}} l l c}
\toprule
密钥层 & 形状 & min \\
\midrule
输入 & 2048, $h{=}42$, 条件间隙 & \textbf{131.8} \\
BSK & 2048, 三元 512, $\sigma{=}2^{-50}$ & \textbf{136.2} \\
packing & 2048, $h{=}256$, $\sigma{=}2^{-44}$ & \textbf{224.3} \\
HWR-KSK & 输入密钥域样本（并入输入层行） & — \\
\bottomrule
\end{tabular*}
\end{table}

\subsection{多参数集修正表（逐环推导）}

$h{=}42$ 是 $n{=}2048$ 输入环族（跨精度）的定案，\textbf{非全局常数}：逐环以三联约束 $T3'\ge128\ \wedge\ \mathrm{MITM}\ge128\ \wedge\ \mathrm{BSK}\ge128$ 求解。$n{=}4096$：原 $h{=}32$ 坍至 $117.5$，修正为 $h{=}42@\mathrm{t}8$（$T3'{=}154.1$；\textbf{已实测}：SQ/686 $=0.93$，噪声 $58.06/58.34$，门测试 $0/4096$ 失配，RS 11 次尝试 $r\_prec{=}8$ 与推导一致）；5-bit（$n{=}2048$，同 $h{=}42$）\textbf{已实测}：SQ/686 $=0.893$，噪声 $56.73/55.65$，门测试 $0/2048$；4096 r-lane stock 对照平价（$1.001$--$1.003$）。$n{=}8192$：修正为 $h{=}34@\mathrm{t}10$（$146.6$，实测排队）。

\section{实验矩阵（终版数据）}

\subsection{性能主表（对比基准 = 686 同机、同场次、同新安全性）}

\begin{table}[htbp]
\centering
\caption{标量（$n{=}2048$ 批，3-bit，中位；dell 静机负载 $\approx 2$）}
\label{tab:scalar}
\small
\begin{tabular*}{\columnwidth}{@{\extracolsep{\fill}} l cccc c}
\toprule
配置 & SQ & 686 & SQ/686 & 噪声 SQ/686 \\
\midrule
Final（$h{=}42,\sigma{+}0$） & \textbf{5.74} ms/msg & 6.40 & \textbf{0.897} & 57.60/57.69 \\
A0 原参数（$39,+0$） & 5.36 & 6.00 & 0.894 & 57.45/57.43 \\
CAL（$41,+11$） & 5.68 & 6.32 & 0.898 & 57.73/58.72 \\
\bottomrule
\end{tabular*}
\end{table}

\begin{table}[htbp]
\centering
\caption{r-lane（$r{=}4$，Final）：同一盲旋转服务 $2048\times4$ 查询。E-PF 实测表明同一算法跨构建二进制的绝对值漂移达 $23\%$，故仅\textbf{同构建内部对照}有效}
\label{tab:rlane}
\small
\begin{tabular*}{\columnwidth}{@{\extracolsep{\fill}} l c c}
\toprule
方案（构建二进制） & ms/msg/lut & 同构建内部对照 \\
\midrule
stock r-lane（stock 构建） & 5.25 & 基准 \\
SQ r-lane（stock 构建） & 5.29 & $1.005\times$（平价） \\
路径A（pathA 构建） & 5.00 & 同构建 stock $4.31$ $\Rightarrow$ $\mathbf{1.16\times}$（更慢） \\
\bottomrule
\end{tabular*}
\end{table}

同场次对 $4\times$686 标量（$52.4$ s，摊销收益）：路径A $1.28\times$，SQ/stock r-lane $1.21\times$，SQ 标量 $\times4$ 为 $1.11\times$。\textbf{安全修正的完整代价}（同机同场）：$h$ 由 39 增至 42 仅 $+6.5\%$（stock）/$+7.2\%$（SQ），$\sigma$ 不动。

\subsection{r-扩展性、E-PF 与正确性}

\textbf{同构建内部对照}（路径A/同构建 stock；r=8 为静机复测 t1 干净批次）：$r{=}2$: $1.12$；$r{=}4$: $1.16$；$r{=}6$: $1.10$；$r{=}8$: $1.05$——路径A 在实测 $r\le 8$ 范围内\textbf{未超过} stock r-lane，差距随 $r$ 收窄，外推交叉点 $r\approx 9$--$12$（未实测验证）。SQ r-lane 对 stock 在全 $r$ 平价（$0.98$--$1.02$，跨三种构建稳定）。每 lane 成本随 $r$ 单调上升（stock 构建 stock：$5.25\to 6.59$ ms/msg，$r{=}4\to 8$），摊销甜点 $r{=}2$--$4$；全部 $r$ 门测试通过（正确性与 $r$ 无关，r=8 亦 Pass）。\textbf{E-PF A/B}（路径A 侧旗标/对象对 SQ 臂的影响）：matflags 构建 SQ 臂 $+0.2\%$--$1.2\%$，stock 构建 $-4.8\%$--$+0.5\%$，而\textbf{同一 stock 算法}跨构建漂移 $-7\%$--$23\%$——证明跨构建绝对值比较无效，且旗标本身近乎免费。

\subsection{自举后容量（余量作为输出维度）}

三层口径：瞬时余量 $= 60-\log_2(\text{最大相位偏差})$；生命周期 $H_{\mathrm{bits}}=\tfrac12\log_2(V_{\mathrm{cap}}/V_{\mathrm{out}})$；容量曲线（自举后施加 $k$ 次乘一 TRGSW 外积逐次测噪）。电路键 $\sigma_G{=}2^{-25}$ 下 8 步噪声平坦（斜率 $<0.05$ bit/EP，外推深度 $\approx 5\times10^4$ 次 EP）；两个零容量边界案例（SAB 级 gadget 单数位截断；输入键 $\sigma{=}2^{-15}$ 放大至半环面）与外积噪声公式 $\sqrt{\ell N}\,B_g/\sqrt{12}\,\sigma_G$ 的饱和点定量验证（$2^{-1.8}\approx$ 观测 $2^{63}$）。\textbf{交替翻转分诊（终参下，5 变体）}：奇数步恒读 $63.00$、偶数步平坦 $57.60$--$57.75$ 的交替模式在 $L{=}8/B_g{=}2^8$ 与 $L{=}16/B_g{=}2^4$ 之间不变，且对 $\sigma_G$ 不敏感直至 $2^{-35}$（此时真实噪声回绕在物理上不可能）——判为\textbf{奇数步测量路径伪象}（符号/参考位问题），非噪声增长；$\sigma_G{=}2^{-15}$ 时全部步饱和（真饱和案例，与公式一致）。容量结论以偶数步链为准，维持不变。

\subsection{与 TFHE-rs 的同机参照与 occupancy}

同机（dell）TFHE-rs v1.8.0 默认 128-bit 经典参数：\textbf{9.560 ms/PBS}（30 次操作精确计数 330 次 PBS，解密校验通过；其参数未按 2026 攻击修正——已标注）。满占用吞吐：路径A $5.00$ ms/查询 $\Rightarrow 1.91\times$；SQ r-lane $5.29\Rightarrow 1.81\times$；SQ 标量 $5.74\Rightarrow 1.67\times$。临界占用率 $\kappa^\ast = T_{\mathrm{batch}}/(nrT_{\mathrm{pbs}})$（$nrT_{\mathrm{pbs}}{=}78.3$ s，$r{=}4$）：路径A $52.3\%$，SQ/stock r-lane $\approx 55\%$；$r{=}1$（SQ 标量）$60.1\%$。即批摊销相对逐消息 PBS 需 $>\!50\%$ 占用率才占优——两种制度（批量吞吐 vs 低延迟）的取舍已如实标注。

\subsection{失效概率（三层证据）}

解析层：实测最大相位偏差 $57.60$ 对预算 $60$，高斯尾部给出逐槽 $p_{\mathrm{fail}}\lesssim 2^{-100}$ 量级（certified 版待定稿数值化）。实证层：累计 $>10^6$ 次槽解码零失败，Clopper--Pearson 单侧 95\% 上界 $p<3\times10^{-6}\approx 2^{-18.4}$（如实标注实证上界与解析主张的区别）。稀有事件估计列为可选第三层。

\section{否定性结果（对社区有价值）}

SQ 与 7 个融合旗标中作用于分解生命周期的部分\textbf{构建不兼容}（同一输出路径冲突，可复现的冲突矩阵）；$k$-sub 推广被蝶形拓扑排除（证明而非性能猜测）；数位宽度微变体已被 stage314 实测关闭（微核 $1.045\times$ 稀释至全 SAB $0.997\times$）；配对堆叠 EP 因成本模型修正而撤回；\textbf{路径A 在终参与 $r\le 8$ 未超越同构建 stock r-lane}（$+5\%$--$16\%$，随 $r$ 收窄）——早前 $1.21\times$ 加速结论系跨构建列误标，已由 E-PF A/B 与通道审计更正（勘误记录于 \S1）。

\section{工件与复现索引}

安全：\texttt{scripts/e1\_conditional\_entropy.py}、\texttt{scripts/e4\_v5*.sage}、dell \texttt{e4v5*.log}/\texttt{e3.log}。性能：\texttt{repro/dell\_phase2/}（主表）、\texttt{repro/stage358..367}（σ 梯度、r-扩展、噪声、容量、平衡档）。理论：\texttt{docs/paper\_ccf\_a/theory\_for\_review.md}、\texttt{review\_package.md}。实现：\texttt{src/sab\_sq.c}、\texttt{src/sab\_pvw\_sq.c}、\texttt{src/sab\_pvw.c}、\texttt{src/mosfhet/src/mattrgsw.c}。全部阶段有 git 提交与服务器日志双重溯源。

\end{document}
